\documentclass[12pt,a4paper]{article}
\usepackage[utf8]{inputenc}
\usepackage[T1]{fontenc}
\usepackage[brazil]{babel}
\usepackage{indentfirst}
\usepackage[margin=2.5cm]{geometry}

\title{Resenha Cr\'itica: Artigo 16 ---\\
ThoughtWorks Technology Radar --- Vol.\ 32}
\author{Gustavo Pessoa Firmino Duarte}
\date{23 de Abril de 2026}

\begin{document}
\maketitle

\section{Introdu\c{c}\~ao}

O \textit{Technology Radar} da ThoughtWorks \'e uma publica\c{c}\~ao semestral
que sintetiza a perspectiva coletiva dos especialistas da empresa sobre tend\^encias
e mudan\c{c}as no mundo da tecnologia de software. O volume 32, publicado em
portugu\^es, organiza as tecnologias em quatro quadrantes --- T\'ecnicas,
Plataformas, Ferramentas e Linguagens \& Frameworks --- e as classifica em quatro
an\'eis: \textit{Adote} (maturidade comprovada, recomenda\c{c}\~ao ativa),
\textit{Experimente} (promissor, vale investigar), \textit{Avalie} (potencial,
mas requer mais observa\c{c}\~ao) e \textit{Mantenha} (n\~ao adote sem raz\~ao
espec\'ifica). O Radar n\~ao \'e um ranking ou uma lista de modismos: \'e uma
destila\c{c}\~ao do julgamento pr\'atico de profissionais que aplicam essas
tecnologias em projetos reais.

\section{Destaques do Volume 32 e o M\'etodo do Radar}

Entre os destaques t\'ecnicos do volume 32, a abordagem de \textit{platform
engineering} --- a constru\c{c}\~ao de plataformas internas que abstraem
complexidades de infraestrutura para os times de produto --- aparece fortemente
posicionada no anel \textit{Adote}, refletindo a maturidade de uma tend\^encia
que come\c{c}ou a surgir em volumes anteriores. A pr\'atica de \textit{architecture
fitness functions}, derivada do livro \textit{Building Evolutionary Architectures}
(Ford et al.), \'e destacada como uma forma de garantir que decis\~oes arquiteturais
n\~ao sejam erodidas ao longo do tempo --- conectando-se diretamente \`a
problem\'atica de d\'eriva arquitetural discutida no artigo de Garlan.

No quadrante de Ferramentas, a consolida\c{c}\~ao de ferramentas de observabilidade
(rastreamento distribu\'ido, m\'etricas e logs unificados) e o crescimento de
plataformas de IA generativa como ferramentas de desenvolvimento aparecem como
tend\^encias dominantes. O Radar tamb\'em registra esc\'epticismo em rela\c{c}\~ao
a certas hypes: algumas ferramentas amplamente discutidas na comunidade aparecem
no anel \textit{Mantenha} ou recebem ressalvas expl\'icitas, demonstrando o valor
do julgamento cr\'itico que o Radar procura cultivar.

O m\'etodo em si \'e instrutivo: o Radar \'e constru\'ido por meio de discuss\~oes
coletivas entre centenas de consultores, buscando deliberadamente pontos de
discord\^ancia e evitando o vi\'es de confirm\'acao de quem trabalha com uma
tecnologia espec\'ifica. Esse processo reflete uma epistemologia madura sobre como
avaliar tecnologias: n\~ao por popularidade ou marketing, mas por experi\^encia
vivida em contextos variados.

\section{Conex\~ao com a Pr\'atica Profissional}

O ThoughtWorks Technology Radar \'e uma das fontes que utilizo regularmente para
calibrar meu mapa mental do ecossistema tecnol\'ogico. Quando escolhi o \textit{stack}
do meu MVP de e-commerce (Node.js, Prisma, React, SQLite), usei o Radar como uma
das refer\^encias para avaliar a maturidade das ferramentas: Prisma estava no anel
\textit{Adote} em volumes anteriores, o que me deu confian\c{c}a de que n\~ao estava
apostando em uma ferramenta experimental.

A discuss\~ao sobre \textit{platform engineering} ressoa com o que observei na
ArcelorMittal: times de suporte a sistemas internacionais que gastavam parte
significativa do tempo configurando ambientes e resolvendo problemas de infraestrutura
que uma plataforma interna bem constru\'ida poderia ter absttra\'ido. A tend\^encia
identificada pelo Radar reflete uma necessidade real que vi em primeira m\~ao.

A pr\'atica de \textit{architecture fitness functions} tamb\'em me parece
particularmente valiosa: a ideia de automatizar verifica\c{c}\~oes das propriedades
arquiteturais --- como depend\^encias proibidas entre m\'odulos ou tamanhos m\'aximos
de componentes --- \'e uma forma concreta de evitar a d\'eriva arquitetural que
tanto o artigo de Garlan quanto a experi\^encia pr\'atica confirmam como inevit\'avel
sem mecanismos de controle ativos.

\section{Conclus\~ao}

O \textit{Technology Radar} da ThoughtWorks \'e um artefato intelectualmente
honesto em um campo frequentemente dominado por modismos e marketing. Seu valor
reside menos nas tecnologias espec\'ificas que lista e mais no modelo de racioc\'inio
que demonstra: como avaliar tecnologias com base em experi\^encia emp\'irica,
contexto de aplica\c{c}\~ao e trade-offs expl\'icitos. Para um engenheiro de
software em forma\c{c}\~ao, aprender a pensar \textit{como o Radar pensa} --- com
cepticismo calibrado, baseado em evid\^encias e sens\'ivel ao contexto --- \'e
possivelmente mais valioso do que qualquer item espec\'ifico que o Radar recomenda.

\end{document}
